\documentclass[12pt]{amsart}
\usepackage[margin=1in,lmargin=0.75in,rmargin=1.25in]{geometry}
\usepackage{amsmath,latexsym,amssymb,amsthm,graphicx}
\usepackage{mathtools}
\DeclarePairedDelimiter\ceil{\lceil}{\rceil}
\DeclarePairedDelimiter\floor{\lfloor}{\rfloor}
\usepackage{subfigure}
\usepackage{xfrac}
\usepackage{fix-cm}
\usepackage{color}
\usepackage[bookmarks=true,%
    colorlinks=true,%
    linkcolor=blue,%
    citecolor=blue,%
    filecolor=blue,%
    menucolor=blue,%
    urlcolor=blue]{hyperref}
\usepackage{doi}
\usepackage{stmaryrd}
\usepackage{multirow}
\usepackage{float}
\usepackage{tabularx}
\usepackage{enumitem}
\usepackage{pdfsync}
\usepackage{tikz}
\usepackage{tikz-cd}
\usepackage{float}
\usepackage[utf8]{inputenc}
\usepackage[textsize=small]{todonotes}
\newcommand{\DPTtodo}[2][]{\todo[color=green!40,#1]{#2}}
\newcommand{\DMGtodo}[2][]{\todo[color=red!40,#1]{#2}}
\usepackage{hyphenat}

\graphicspath{{draws/}{}}

\setenumerate[1]{label=(\arabic*)}

\numberwithin{equation}{section}
\numberwithin{figure}{section}

\newtheorem{theorem}{Theorem}[section]
\newtheorem{lemma}[theorem]{Lemma}
\newtheorem*{claim}{Claim}
\newtheorem{corollary}[theorem]{Corollary}
\newtheorem{proposition}[theorem]{Proposition}
\newtheorem{exercise}[theorem]{Exercise}
\newtheorem*{ex*}{Exercise}

\theoremstyle{definition}
\newtheorem{remark}[theorem]{Remark}
\newtheorem{remarks}[theorem]{Remarks}
\newtheorem{example}[theorem]{Example}
\newtheorem{definition}[theorem]{Definition}
\newtheorem{question}[theorem]{Question}
\newtheorem{conjecture}[theorem]{Conjecture}
\newtheorem{notation}[theorem]{Notation}
\newtheorem{fact}[theorem]{Fact}

\newcommand\N{\ensuremath{\mathbb{N}}}
\newcommand\Z{\ensuremath{\mathbb{Z}}}
\newcommand\Q{\ensuremath{\mathbb{Q}}}
\newcommand\R{\ensuremath{\mathbb{R}}}
\newcommand\C{\ensuremath{\mathbb{C}}}
\def\H{\ensuremath{\mathbb{H}}} % Redefining "long Hungarian umlaut"
\newcommand\D{\ensuremath{\mathbb{D}}}
\newcommand\F{\ensuremath{\mathcal{F}}}
\newcommand\G{\ensuremath{\mathcal{G}}}
\newcommand\J{\ensuremath{\mathcal{J}}}
\newcommand\U{\ensuremath{\mathcal{U}}}
\newcommand\V{\ensuremath{\mathcal{V}}}
\DeclareMathOperator{\GL}{GL}
\DeclareMathOperator{\Sp}{Sp}
\DeclareMathOperator{\SL}{SL}
\DeclareMathOperator{\Curr}{Curr}
\DeclareMathOperator{\Teich}{Teich}
\newcommand{\pa}[1]{\left(#1\right)}
\newcommand{\cb}[1]{\left\{#1\right\}}
\newcommand{\tn}[1]{\textnormal{#1}}
\newcommand{\br}[1]{\left[#1\right]}
\newcommand{\rest}[1]{\left.#1\right|}
\newcommand{\card}[1]{\left| #1 \right|}
\newcommand{\norm}[1]{\left\| #1 \right\|}
\newcommand{\ab}[1]{\langle #1 \rangle}
\renewcommand{\card}[1]{\tn{card}\pa{#1}}
\renewcommand{\Re}{\tn{Re}}
\renewcommand{\Im}{\tn{Im}}
\newcommand{\ang}{\tn{ang}}
\newcommand{\Arg}{\tn{Arg\,}}
\renewcommand{\arg}{\tn{arg\,}}
\newcommand{\Log}{\tn{Log\,}}
\newcommand{\len}{\tn{len\,}}
\newcommand{\del}{\partial}
\newcommand{\Res}[2]{\tn{Res}\br{#1, #2}}

\renewcommand{\epsilon}{\varepsilon}

\def\mathcenter#1{%
  \vcenter{\hbox{$#1$}}%
}
\def\mfig#1{
        \mathcenter{\includegraphics{#1}}
}

\usetikzlibrary{arrows}
\tikzset{
    labl/.style={anchor=south, rotate=90, inner sep=.5mm}
}
\tikzstyle{every picture}=[> = to]
% Style for labels on arrows in commutative diagrams
\tikzset{cdlabel/.style={execute at begin node=$\scriptstyle,execute at end node=$}}
\tikzset{implication/.style={double equal sign distance, -implies}}
\tikzset{biimplication/.style={double equal sign distance, implies-implies}}

\begin{document}
\title{Application 1: convexity along paths}
\author{}
\maketitle


This contains the conjectures and possible applications related to characterization of convexity of paths in Teichmueller space.

\subsection{Main results}

\DPTtodo{Define ``curve'', ``essential crossing''}
\begin{definition}\label{def:convex}
  Let $f \colon C^+(S) \to \R_{\ge 0}$ be a curve functional. We say that $f$ is \emph{convex}
  if it satisfies the following two properties.
  \begin{itemize}
  \item \textbf{(Smoothing)} Let $C$~be an oriented curve on~$S$, and
    let $x$~be an essential crossing of~$C$. Let $C'$ be the oriented
    smoothing of~$C$ at~$x$. Then $f(C) \ge f(C')$. Schematically, we
    have
    \begin{equation}\label{eq:smoothing}
      f\left(\mfig{curves-3}\right) \ge f\left(\mfig{curves-4}\right)
    \end{equation}
  \item \textbf{(Convex union)} Let $C_1$ and~$C_2$ be two curves
    on~$S$. Then
    \begin{equation}\label{eq:union}
      f(C_1 \cup C_2) \le f(C_1) + f(C_2).
    \end{equation}
  \end{itemize}
\end{definition}

An example of a functional that is not convex is stable word-length w.r.t to a non-simple generating set (example of genus 1).



\begin{definition}
  A homogeneous function on geodesic currents is \emph{convex} if its
  restriction to curves is convex by Definition~\ref{def:convex}.
\end{definition}


\begin{definition}
We say a path $\mu_t$ in the space $\mathcal{G}\mathcal{C}(\Sigma)$ is \emph{convex} if for every convex functional on $\mathcal{G}\mathcal{C}(\Sigma)$, $t \mapsto f(\mu_t)$ is convex.
\end{definition}

We claim it is enough to look at intersection numbers with closed curves, i.e.

\begin{lemma}
For all $t \mapsto f(\mu_t)$ is convex for all convex functional $f$ if and only if $t \mapsto i(\mu_t,\alpha)$ is convex for all closed geodesic $\alpha$.
\end{lemma}

A way to rephrase this is as follows

\begin{lemma}
 If $ i(\mu_{\frac{s+t}{2}},\alpha) \leq \frac{1}{2} i(\mu_{s}+\mu_{t},\alpha) $ then $ f(\mu_{\frac{s+t}{2}}) \leq \frac{1}{2} f(\mu_{s}+\mu_{t})$ for all convex functional $f$.
\end{lemma}

Another way to rephrase it is

\begin{lemma}
 Let $\mu_1,\mu_2,\mu_3$ geodesic currents, and $\alpha$ any geodesic.
 If $ i(\mu_1, \alpha) \leq \frac{1}{2}(i(\mu_2, \alpha)+i(\mu_3, \alpha))$, then $f(\mu_1) \leq \frac{1}{2} (f(\mu_2)+f(\mu_3))$ for all convex functional $f$.
\end{lemma}

Fixing a hyperbolic metric $m$ on $\Sigma$, we get the function hyperbolic length $L_m\colon \mathbb{R}\mathcal{S}'(\Sigma) \to \mathbb{R}$.
We can also consider the function extremal length
$EL_m\colon \mathbb{R}\mathcal{S}'(\Sigma) \to \mathbb{R}$.
By a result by Lenzhen-Rafi we have that there exist some Teichm\"uller geodesic $t \mapsto m_t$ in Teichmueller space, and a simple curve $\alpha$, such that  $t \mapsto L_{m_t}(\alpha)$  and  $t \mapsto EL_{m_t}(\alpha)$  are not convex. \\
Recall that $L_{m_t}(\alpha)=i(m_t,\alpha)$. According to our previous definitions, this says that a path of Liouville currents is not convex (i.e. Teichm\"uller geodesics are not convex).

On the other hand, it is known that hyperbolic length is convex w.r.t to earthquake paths (Kerchoff), WP-geodesics (Wolpert) and shearing paths (Bestvina-Bromberg-Fujiwara-Souto).

\section{Non-convexity of Teichm\"uller geodesics}



\section{Convexity of Earthquake paths}

By fixing a pants decomposition with a simple closed curve $C$ as a pants boundary, we can define the earth-quake of $g \in \mathcal{T}(\Sigma)$ 
$\mathcal{E}_{C^t}(g)$ be the hyperbolic metric obtained by twisting a distance $t$ (measured in $g$) along $C$.
Kerchoff proves convexity of $l_{\gamma}\colon \mathcal{T}(\Sigma) \to \mathbb{R}$ by studying its first order variation and seeing it's a strictly increasing function. \\
Bestvina-Bromberg-Fujiwara-Souto prove also this in the more general case of measured laminations with finitely many leaves.

\section{Stretch lines or stretch-geodesics}



\section{Convexity of WP-geodesics}

(Wolpert)

\section{Convexity of shearing paths (can we subsume here earth-quake paths and stretch-lines?)}


\begin{remark}
This characterization of convexity does not a priori shed light to the question is extremal length convex along the above paths.
\end{remark}

\section{Convexity on $ML$ w.r.t to natural coordinates}

\end{document}


%\end{thebibliography}
\end{document}
